\section{Ordered bar coalgebras and Lie chains}
\label{chap:ordered-lie-bar}

The enveloping coproduct identifies the Lie generators as primitives.
The ordered associative bar has another coproduct, obtained by cutting
a tensor into a prefix and a suffix. We construct its differential and
the explicit map from Lie chains, including the signs and factorials.
The map is a coalgebra quasi-isomorphism, but it does not identify the
strict primitive complexes. This is a classical comparison of an
enveloping algebra with its ordered bar; it does not identify either
construction with an associative chiral duality.

\subsection{Products transported through the ordered basis}
\label{sec:olb-transported-product}

Let $\mathfrak g$ be a differential graded Lie algebra over $\kk$.
The symmetrization map of
Proposition~\ref{sp:symmetrization-theorem} is an isomorphism of
coalgebra complexes
$\sigma:S(\mathfrak g)\to U(\mathfrak g)$.
To retain the associative product as well, define
\begin{equation}
 a\star b=\sigma^{-1}\bigl(\sigma(a)\sigma(b)\bigr).
 \label{eq:olb-transported-product}
\end{equation}
This is an actual operation on the symmetric vector space.
It depends on the Lie bracket through the ordered relations.

\begin{proposition}\label{prop:olb-transported-product}
With this product, $S(\mathfrak g)$ is a differential graded Hopf
algebra and $\sigma$ is a Hopf isomorphism.
Its unit, counit, differential and coproduct are the ones already
defined on $S(\mathfrak g)$.
The product respects the increasing symmetric-length filtration.
For homogeneous $x,y\in\mathfrak g$,
\[
 x\star y=xy+\tfrac12[x,y].
\]
\end{proposition}

\begin{proof}
Applying $\sigma$ to $(a\star b)\star c$ and $a\star(b\star c)$
gives the same product $\sigma(a)\sigma(b)\sigma(c)$.
Injectivity proves associativity.  The same calculation with the
unit proves its two identities.
Since $\sigma$ commutes with the differential, applying it to the
graded Leibniz equation for $\star$ gives that equation in
$U(\mathfrak g)$.

The coalgebra identity for $\sigma$ transports bialgebra
compatibility, and its counit remains multiplicative.
Transporting the antipode gives the two antipode equations.
Both $\sigma$ and its inverse preserve the filtration, so the
product has the stated bound.
Finally,
\[
 \sigma(xy)=\tfrac12\bigl(xy+(-1)^{|x||y|}yx\bigr)
            =xy-\tfrac12[x,y]
\]
in the enveloping algebra.  Since $\sigma$ fixes each Lie element,
this proves the displayed formula.
\end{proof}

A graded vector-space identification therefore need not preserve
the product that enters an iterated multiplication.
The half-bracket is already visible on two inputs.
It will also be visible in the differential of an ordered bar word.

\subsection{Augmentation and the ordered bar differential}
\label{sec:olb-ordered-bar}

Let $A$ be a unital differential graded associative algebra over $\kk$.
An augmentation is a unital algebra chain map
$\varepsilon:A\to \kk$, where $\kk$ has degree zero and zero differential.
Its kernel $\bar A$ is an ideal preserved by the differential.
The splitting by the unit gives $A=\kk1\oplus\bar A$ as complexes.
Write $sa$ for the downward shift of $a\in\bar A$,
so $|sa|=|a|-1$ and the shift differential is $d(sa)=-s(da)$.

Define the ordered bar coalgebra
\[
 B_-(A)=\kk\oplus\bigoplus_{r\geq1}(s\bar A)^{\otimes r}.
\]
The subscript records the minus sign in the shifted differential.
All elements have finite word support and all tensors are algebraic.
Write $[a_1|\cdots|a_r]=sa_1\otimes\cdots\otimes sa_r$.

The empty word is its coaugmentation and its counit kills nonempty
words.  Deconcatenation cuts every word at each gap:
\begin{equation}
 \Delta[a_1|\cdots|a_r]
 =\sum_{j=0}^r[a_1|\cdots|a_j]\otimes[a_{j+1}|\cdots|a_r].
 \label{eq:olb-deconcatenation}
\end{equation}
In \eqref{eq:olb-deconcatenation}, two successive cuts amount to
the same ordered pair of gaps in either order, proving coassociativity.

Define degree-one local operations by
\[
 b_1(sa)=-s(da),\qquad
 b_2(sa\otimes sb)=(-1)^{|a|}s(ab).
\]
Insert either operation after a prefix of shifted degree $q$ with
the sign $(-1)^q$.
In particular,
\begin{equation}
 b_{\mathrm{mult}}[a_1|\cdots|a_r]
 =\sum_{i=1}^{r-1}
 (-1)^{\sum_{j<i}(|a_j|-1)+|a_i|}
 [a_1|\cdots|a_ia_{i+1}|\cdots|a_r].
 \label{eq:olb-bar-faces}
\end{equation}
Multiplicativity of the augmentation ensures that every merged
output still belongs to $\bar A$.

\begin{proposition}\label{prop:olb-bar}
The sum $b_-=b_1+b_{\mathrm{mult}}$, zero on the empty word,
makes $B_-(A)$ a differential graded coalgebra.
An augmented algebra chain map $f:A\to A'$ induces a coalgebra
chain map by applying $f$ to every tensor factor.
These maps preserve identities and compositions.
\end{proposition}

\begin{proof}
Two internal differentials on the same factor give zero.
Two operations on disjoint intervals cancel in the square because
both have odd degree.
For three overlapping inputs, the multiplication contribution is
\[
 b_{\mathrm{mult}}^2[a|b|c]
 =(-1)^{|b|}s\bigl((ab)c-a(bc)\bigr)=0.
\]

For example, when $a,b,c$ all have odd degree,
\[
 b_{\mathrm{mult}}[a|b|c]=-[ab|c]-[a|bc],\qquad
 b_{\mathrm{mult}}^2[a|b|c]=-[(ab)c]+[a(bc)]=0.
\]
Here $ab$ has even degree, which gives the opposite signs in the
second equality.

For the mixed operations on a pair, put $p=|a|$. The Leibniz rule gives
\begin{align*}
 b_1b_2(sa\otimes sb)&=(-1)^{p+1}s((da)b)-s(a\,db),\\
 b_2b_1(sa\otimes sb)&=(-1)^ps((da)b)+s(a\,db).
\end{align*}
Their sum is zero.
These cases account for every term of $b_-^2$ using
\eqref{eq:olb-bar-faces}.

A cut outside a local operation survives before or after that
operation with the same sign.
A cut inside a merged pair survives in neither expression.
Thus $\Delta b_-=(b_-\otimes1+1\otimes b_-)\Delta$.
Applying $f$ to tensor factors commutes with $b_1$, with multiplication,
and with every cut.  The same formulas prove the assertions about
identities and compositions.
\end{proof}

This ordered coproduct differs from the primitive coproduct of an
enveloping algebra.  In fact,
\begin{equation}
 P(B_-(A))=s\bar A
 \label{eq:olb-bar-primitives}
\end{equation}
as complexes.  To prove it, project the reduced coproduct of a
length-$r$ word to $s\bar A\otimes(s\bar A)^{\otimes(r-1)}$.
For $r\geq2$ this projection is the original word with its first
factor separated, hence is injective.
Different lengths have different total output lengths.
Only length one remains primitive, with differential $-sd$.

For $A=\kk[x]$, with $|x|=0$ and zero differential,
$[x^2]$ is a nonzero class in $H^{-1}(P(B_-(A)))$.
Nevertheless,
\[
 b_-[x|x]=[x^2],\qquad
 \bar\Delta[x|x]=[x]\otimes[x]\ne0.
\]
Thus the map from cohomology of strict bar primitives to primitives
in full bar cohomology is not injective.
Theorem~\ref{sp:cohomology} applies to the enveloping Hopf
coproduct. It does not imply the corresponding assertion for the
deconcatenation coproduct of the bar.

Applying Proposition~\ref{prop:olb-bar} to the product
\eqref{eq:olb-transported-product} gives an actual chain isomorphism
\[
 B_-(S(\mathfrak g),\star)\xrightarrow{B_-(\sigma)}B_-(U(\mathfrak g)).
\]
If the source instead uses the ordinary symmetric product, its
chain defect in length two on $[x|y]$ is
\[
 (-1)^{|x|}s\bigl(\sigma(x)\sigma(y)-\sigma(xy)\bigr)
 =\tfrac12(-1)^{|x|}s[x,y].
\]
This computes the missing term without identifying two different
bar differentials.

\subsection{The sign convention for higher operations}
\label{sec:olb-sign-convention}

The positive-length part of an ordered tensor coalgebra has no
counit.  Its reduced coproduct cuts only into two nonempty words.
For a graded space $V$, write it as
$T^c(sV)=\bigoplus_{r\geq1}(sV)^{\otimes r}$.
A degree-one map on this coalgebra is a coderivation if
\[
 \bar\Delta b=(b\otimes1+1\otimes b)\bar\Delta.
\]
A map from $(sV)^{\otimes n}$ to $sV$ extends to a coderivation by
inserting it into every consecutive interval, with its degree-one
prefix sign.  Cuts before and after that interval prove the
coderivation identity.

Conversely these components determine a coderivation.
Indeed, the difference on a shortest word where two such maps
disagree has zero coproduct.  It is therefore a word of length one,
by the first-cut argument proving \eqref{eq:olb-bar-primitives}.
Its projection to length one is zero, a contradiction.

An uncurved $A_\infty$-structure on $V$ is a degree-one coderivation
on $T^c(sV)$ with square zero. There is no zero-input operation in
this definition. Define its unshifted operations by
\[
 b_n s^{\otimes n}=s m_n,\qquad |m_n|=2-n.
\]
For homogeneous maps, the convention is
$(F\otimes G)(v\otimes w)=(-1)^{|G||v|}F(v)\otimes G(w)$.
It gives
\[
 s^{\otimes n}(x_1\otimes\cdots\otimes x_n)
 =(-1)^{\sum_j(n-j)|x_j|}sx_1\otimes\cdots\otimes sx_n.
\]

For a differential graded algebra with operations $m_1=d$ and
$m_2=\mu$, this convention has
$b_1(sa)=+s(da)$ and the same binary operation
$b_2(sa,sb)=(-1)^{|a|}s(ab)$.
Call this differential $b_+$.
Its square vanishes: changing the unary sign multiplies every
mixed cancellation above by minus one and preserves the other
zero sums.

If the algebra is augmented, $T^c(s\bar A)$ is a subcoalgebra
preserved by this differential, since $d$ and products preserve
$\bar A$.  Adjoining the empty word gives a counital coalgebra
$B_+(A)$.  The algebra unit as an input of $T^c(sA)$ is not that
empty word, and is excluded on restriction to $s\bar A$.

For a degree-zero coalgebra map
$\Phi:T^c(sV)\to T^c(sW)$, let $\Phi_n$ be its projection to
one output factor on $n$ inputs.  Define its unshifted component by
$\Phi_n s^{\otimes n}=s f_n$.  It has degree $1-n$.

\begin{proposition}[Changing the unary sign]
\label{prop:olb-sign-change}
The formula
\[
 R[a_1|\cdots|a_r]=(-1)^{\sum_i|a_i|}[a_1|\cdots|a_r],
 \qquad R(1)=1,
\]
defines a degree-zero coalgebra involution and a chain isomorphism
$R:B_-(A)\to B_+(A)$.
For general higher operations on $T^c(sV)$, conjugation by $R$
changes $m_n$ to $(-1)^n m_n$.
For a coalgebra morphism with unshifted component $f_n$ of degree
$1-n$, conjugating its source and target changes that component
to $(-1)^{n-1}f_n$.
\end{proposition}

\begin{proof}
The exponent splits additively at each cut, so $R$ preserves the
coproduct.  It squares to the identity.
An internal differential increases the sum of unshifted degrees
by one, whereas a multiplication preserves that sum.
Therefore $Rb_-R=b_+$.

More generally, an $n$-ary operation changes the sum by $2-n$.
The signs on its inputs and output have ratio
$(-1)^{2-n}=(-1)^n$.
For a morphism component the same ratio is
$(-1)^{1-n}=(-1)^{n-1}$.
Conjugation preserves the square-zero and chain-map equations,
so these formulas hold at every arity.
\end{proof}

For $da=b$, with $|a|=0$ and $|b|=1$, one has
$b_-[a]=-[b]$, $R[a]=[a]$, and $R[b]=-[b]$.
Thus both sides of $Rb_-[a]=b_+R[a]$ equal $[b]$.
In higher arity, a fourth or sixth operation retains its sign,
while a second comparison component changes sign.
The translation must therefore include the comparison maps, not
only the multiplication.

\subsection{Relative coefficients and parity base change}
\label{sec:olb-relative-signs}

Let $C$ be a commutative differential graded $\kk$-algebra.
Let $A$ be a differential graded algebra with graded central
$C$-action and an augmentation $A\to C$ splitting its unit.
Thus $ac=(-1)^{|a||c|}ca$ for homogeneous elements and
$A=C\oplus\bar A$ as $C$-complexes.

The shifted scalar action is
$s(ca)=(-1)^{|c|}c\,sa$.

The relative bar is
\[
 B_C(A)=C\oplus\bigoplus_{r\geq1}(s\bar A)^{\otimes_C r}.
\]
It has the coefficient differential on length zero and the
operations $-sd$ and $(-1)^{|a|}s(ab)$ on shifted inputs.
The proof of Proposition~\ref{prop:olb-bar} uses the coefficient
Leibniz rule as well, and applies term by term over $C$.
Its coproduct takes values in a tensor product over $C$.

Put $C^-=(C,-d_C)$ and $A^-=(A,-d_A)$, retaining the products
and gradings.  The sign change preserves the Leibniz equation.
Define parity maps
\[
 \pi_C(c)=(-1)^{|c|}c,\qquad \pi_A(a)=(-1)^{|a|}a.
\]
They are algebra isomorphisms, with
$(-d_C)\pi_C=\pi_Cd_C$ and
$(-d_A)\pi_A=\pi_Ad_A$.
They also preserve the augmentation through the base map $\pi_C$.

\begin{proposition}\label{prop:olb-relative-parity}
Applying these parity maps to coefficients and every tensor factor gives
an isomorphism of relative coalgebra complexes
\[
 B_C(A)\longrightarrow B_{C^-}(A^-).
\]
The map is semilinear over $\pi_C$: it sends $cv$ to
$\pi_C(c)$ times the image of $v$.
The target unary operation is $+s d_A$, and its length-zero
differential is $-d_C$.
\end{proposition}

\begin{proof}
The parity map intertwines the module actions by
$\pi_A(ca)=\pi_C(c)\pi_A(a)$.
It therefore defines maps on balanced tensor products over the
displayed base isomorphism.
Compatibility with the shifted scalar action follows from
$s(ca)=(-1)^{|c|}c\,sa$ on both sides.
The algebra and chain identities for $\pi_A$ show that the maps
commute with the two local bar operations.
They commute with cuts because they act on each tensor factor.
The same maps applied again give the inverse.

Finally the standard target unary sign is
$-s(-d_A)=+s d_A$, while its coefficient differential is $-d_C$.
\end{proof}

Changing only the unary formula while silently fixing an odd
coefficient is not this base change.
For example, a degree-one coefficient $\omega$ is sent to
$-\omega$, and a map to a field concentrated in degree zero
must kill $\omega$ altogether.
Thus a field-valued specialization cannot retain a nonzero
degree-one collision coefficient.
The relative bar keeps that coefficient in its own differential
graded base algebra and its tensor products over that base.

A map of commutative differential graded algebras $C\to E$
has a different, covariant bar operation.
Writing $A_E=E\otimes_C A$, the split augmentation gives
$\overline{A_E}=E\otimes_C\bar A$.
The balanced tensor identifications then define
\[
 E\otimes_C B_C(A)\xrightarrow{\sim}B_E(A_E).
\]
They send each word to the same word with extended coefficients.
The differential formulas and every cut agree, so this is an
isomorphism of coalgebra complexes.

This is an assertion about the ordinary split tensor construction.
If a derived coefficient extension is intended, it requires an
appropriate resolution before tensoring.

\subsection{Lie chains inside ordered bar words}
\label{sec:olb-lie-chains}

For $A=U(\mathfrak g)$, the augmentation kills every Lie
generator.  The inclusion $\mathfrak g\hookrightarrow\bar A$
therefore gives shifted Lie elements inside its bar coalgebra.
The Lie bracket controls a cocommutative subcoalgebra of these
ordered words.

Put $W=s\mathfrak g$ and give $S(W)$ its primitive-generator
symmetric coproduct.  Define
\begin{equation}
 \chi(w_1\cdots w_r)=
 \sum_{\pi\in\mathfrak S_r}\epsilon(\pi;w)
 w_{\pi(1)}|\cdots|w_{\pi(r)},\qquad \chi(1)=1,
 \label{eq:olb-lie-bar-map}
\end{equation}
where the shifted Lie elements on the right are included in
$s\overline{U(\mathfrak g)}$.
The signs are the Koszul signs of the shifted degrees
$|sx|=|x|-1$.  There is no factorial denominator in this formula.
Reindexing permutations shows that it is well-defined on $S(W)$.
In a repeated shifted element of odd degree the terms cancel, as required
by the symmetric relation on $W$.

Cutting a permuted word separates its labels into a subset and
its complement, each with an internal order.
These are exactly the terms obtained from the symmetric
coproduct and two applications of $\chi$.
The subset and internal permutation signs multiply to the
original permutation sign.  Thus \eqref{eq:olb-lie-bar-map}
is a coalgebra map.
It is injective: the quotient from $W^{\otimes r}$ to $S^r(W)$
composed with its permutation sum is $r!$ times the identity.
The inclusion of Lie elements in the bar is injective by
Theorem~\ref{sp:pbw}, and tensors of vector-space injections
are injective by extension of bases.

Define degree-one maps on one and two symmetric shifted inputs by
\[
 \ell_1(sx)=-s(dx),\qquad
 \ell_2(sx\,sy)=(-1)^{|x|}s[x,y].
\]
Graded skew symmetry of the bracket makes the second map
symmetric with the signs of the shifted inputs.

For $n=1,2$, extend it to $S(W)$ by
\begin{equation}
 D_n(w_1\cdots w_r)=
 \sum_{\substack{J\subseteq\{1,\ldots,r\}\\|J|=n}}
 \epsilon(J,J^c;w)\ell_n(w_J)w_{J^c}.
 \label{eq:olb-lie-chain-differential}
\end{equation}
The entries in each subset keep their original order.
Set the result to zero when $r<n$, and set $D_n(1)=0$.
Reindexing labeled factors proves compatibility with the
symmetric relations.  Each summand has degree one.

The extensions \eqref{eq:olb-lie-chain-differential} are coderivations.  A term of $\Delta D_n$
selects the contracted subset and then separates all remaining
labels into two sets.  If the contracted output lies on the
left, it is the corresponding term of $(D_n\otimes1)\Delta$.
If it lies on the right, moving the degree-one operation past
the left factors gives exactly the tensor sign of
$(1\otimes D_n)\Delta$.
Every term occurs in one of these two cases.

A chain map is a quasi-isomorphism if it induces an isomorphism
in every cohomological degree.

\begin{theorem}[Lie chains and the enveloping bar]
\label{thm:olb-lie-bar}
The operator $D=D_1+D_2$ has square zero on $S(s\mathfrak g)$.
With this differential, the map
\[
 \chi:(S(s\mathfrak g),D)\longrightarrow B_-(U(\mathfrak g))
\]
is a natural injective coalgebra chain map and a quasi-isomorphism.
The source is the Chevalley--Eilenberg chain coalgebra of
$\mathfrak g$ with the displayed suspension convention.
All tensors and sums here are algebraic, and $\mathfrak g$
may be unbounded and infinite-dimensional.
\end{theorem}

\begin{proof}[Proof of the chain identity]
On one shifted Lie element the two differentials agree.
On two inputs, put $p=|x|$ and $q=|y|$.
The length-one part of the bar differential is
\begin{align*}
 b_{\mathrm{mult}}\chi(sx\,sy)
 &=(-1)^p s(xy)
   +(-1)^{(p-1)(q-1)+q}s(yx)\\
 &=(-1)^p s\bigl(xy-(-1)^{pq}yx\bigr)
  =(-1)^p s[x,y].
\end{align*}
It equals the length-one part of $\chi D$.
For more than two inputs both length-one parts vanish.

To see that these checks prove the full equation, put
$E=b_-\chi-\chi D$.
The coderivation identities and the coalgebra identity for
$\chi$ give
\[
 \bar\Delta E=(E\otimes\chi+\chi\otimes E)\bar\Delta.
\]
If $E$ vanishes in all smaller symmetric lengths, its value in
the next length has zero reduced coproduct.
It is then a bar element of length one, by the first-cut argument.
Its projection to length one was just proved to be zero.
Induction proves $E=0$.

Now $\chi D^2=b_-^2\chi=0$, and injectivity of $\chi$ gives
$D^2=0$.  The formulas commute with differential graded Lie
maps, proving naturality.
\end{proof}

The quasi-isomorphism requires more than this chain identity.
An explicit resolution proves it for an abelian graded Lie
algebra, and finite filtrations then add the differential and
the bracket.

Here a resolution of the augmentation module $\kk$ means an
augmented complex $P\to \kk$ inducing a cohomology isomorphism.
A chain homotopy from $G$ to $F$ is a degree-minus-one map $L$
satisfying $dL+Ld=F-G$.  For maps of graded left $A$-modules,
its linearity is $L(av)=(-1)^{|a|}aL(v)$.

\begin{lemma}[The abelian bar comparison]
\label{lem:olb-abelian-bar}
For a graded vector space $V$ with zero differential,
\[
 (S(sV),0)\xrightarrow{\chi}B_-(S(V))
\]
is a quasi-isomorphism.  It preserves total polynomial weight,
where every vector of $V$ has weight one.
\end{lemma}

\begin{proof}
Write $A=S(V)$ and consider the graded commutative algebra
$K=A\otimes S(sV)$.
Give it the degree-one derivation $\partial$ defined by
\[
 \partial x=0,\qquad \partial(sx)=-x.
\]
Its square is zero on the generators, and the mixed Leibniz
terms cancel, so its square is zero everywhere.

Let $h$ be the degree-minus-one derivation defined by
$h(x)=-sx$ and $h(sx)=0$.
The anticommutator $\partial h+h\partial$ is a degree-zero
derivation that fixes each generator.
It consequently acts by multiplication by $N$ on total
polynomial weight $N$.

Thus $h/N$ contracts every positive weight, while weight zero
is $\kk$.  Projection $K\to \kk$ is a resolution of the augmentation
module, free as a graded left $A$-module on the ordered
monomials of $S(sV)$.

There is also an ordered bar resolution
$P=A\otimes B_-(A)$.
Its differential consists of the bar multiplication on the
second factor, with the tensor sign, and the leading face
\[
 a_0\otimes[a_1|\cdots|a_r]
 \longmapsto-(-1)^{|a_0|}a_0a_1\otimes[a_2|\cdots|a_r].
\]
The leading face is zero in length zero.
Its square cancels with the first internal bar face by
associativity.  Disjoint faces cancel by their odd degrees.
The augmentation $P\to \kk$ is the usual augmentation on
length zero and is zero on nonempty bar words.

Put $\bar a_0=a_0-\varepsilon(a_0)1$.
The degree-minus-one map
\[
 H_P(a_0\otimes[a_1|\cdots|a_r])
 =-1\otimes[\bar a_0|a_1|\cdots|a_r]
\]
contracts this augmentation as a $\kk$-complex.
The leading face contributes $\bar a_0\otimes[a_1|\cdots|a_r]$.
All faces strictly to its right cancel in the two compositions.

When $r>0$, the next face contributes the remaining
$\varepsilon(a_0)\otimes[a_1|\cdots|a_r]$.
Indeed the coefficient can be nonzero only in degree zero,
and its two appearances agree by the augmentation splitting.
In length zero the result is $a_0-\varepsilon(a_0)1$.
Thus $d_PH_P+H_Pd_P$ is the identity minus the projection to
the constant summand.

Both this contraction and $h/N$ preserve polynomial weight
and the sum of unshifted degrees.

Define $i:K\to P$ by $1_A\otimes\chi$.
It preserves the augmentation.  Its internal bar contribution
vanishes because $V$ is abelian, by the chain identity already
proved.  Selecting the first tensor factor of a permuted word gives
the terms of $\partial$ with their graded Leibniz signs.
Hence $i$ is an $A$-linear chain map.

Here is an explicit comparison argument that permits tensoring
these resolutions with $\kk$.
Give their free generators an auxiliary resolution length:
the number of shifted factors in $K$ and the bar length in $P$.
Both differentials decrease this nonnegative length by one.
In length zero both modules are $A$, and take the comparison
inverse $j:P\to K$ to be the identity there.

If $j$ is defined on all shorter free generators, then
$j(d_Pw)$ is a cycle with zero augmentation.
Set $j(w)$ equal to its image under the contraction $h/N$
in its polynomial weight $N>0$.
It satisfies $\partial j(w)=j(d_Pw)$.
Extend over $A$ and repeat in increasing resolution length.
Every value is a finite expression.

The same induction constructs homotopies from $ji$ and $ij$
to the respective identities.
For two augmentation-preserving chain maps $F,G$ and a free
generator $w$, suppose the homotopy $L$ has been defined in
smaller lengths.
The element
\[
 z=(F-G)(w)-L(dw)
\]
is a cycle with zero augmentation, because the homotopy
equation already holds on $dw$.
Apply the target contraction to $z$ to define $L(w)$.

Then $dL(w)+L(dw)=(F-G)(w)$.
Extend this degree-minus-one map by
$L(aw)=(-1)^{|a|}aL(w)$.
This verifies the homotopy equation on all free-module elements.
The construction respects polynomial weight and unshifted degree.

Thus $i$ is an $A$-linear homotopy equivalence.
Tensor it over $A$ with the right augmentation module $\kk$.
The resulting maps and homotopies remain defined and retain
their equations.  The resulting complexes are $S(sV)$ with
zero differential and $B_-(A)$ with its bar differential.
The resulting map is exactly $\chi$, proving the lemma.
\end{proof}

\begin{proof}[Proof of the quasi-isomorphism in
Theorem~\ref{thm:olb-lie-bar}]
First let $\mathfrak g=V$ have zero bracket but a possibly
nonzero differential.  Fix total polynomial weight $N$.
The source has $N$ shifted factors.  Every target bar word has
at most $N$ factors, because each augmentation-ideal factor
has positive polynomial weight.
Separate the internal differential from the bar multiplication.
They anticommute by the mixed part of the bar square calculation.

The cone of $\chi$ has two corresponding differentials.
One consists of multiplication faces and the comparison map,
and the other is the internal differential.
Here the cochain cone of $f:X\to Y$ is
$\operatorname{Cone}(f)^n=Y^n\oplus X^{n+1}$, with
$d(y,x)=(d_Yy+f(x),-d_Xx)$.
The first differential preserves the sum $p$ of unshifted
degrees.  In each such degree its complex is acyclic by
Lemma~\ref{lem:olb-abelian-bar}.

The internal differential increases $p$ by one.
In total cochain degree $n$, the range of $p$ is finite.
The bar length lies between zero and $N$, and the cone adds
at most one additional shift.

For a cycle, take its smallest $p$ component.
It is a cycle for the first differential.
Choose a primitive for it using that acyclic complex and
subtract its full differential.
This removes the smallest component, and the extra internal
differential has strictly larger $p$.

The process terminates because the possible $p$ values in
degree $n$ form a finite interval.
It expresses every cone cycle as a boundary, proving the
comparison for an abelian differential graded Lie algebra.
This finite elimination works even when $V$ is unbounded.

For a nonzero bracket, filter $S(s\mathfrak g)$ by symmetric
length.  Filter the bar target by the sum of the word
filtrations of its enveloping-algebra factors.
Each factor in the augmentation ideal has filtration at
least one.  The differential preserves these filtrations,
and $D_2$ lowers the source filtration by one.

Theorem~\ref{sp:pbw} identifies the associated graded target
with the bar of $S(\mathfrak g)$, retaining the internal
differential.  This follows on tensor powers from the ordered
bases, with the sum filtration on their factors.
The associated graded comparison is precisely the abelian
differential graded comparison just proved.

A cycle in the cone has finite greatest filtration degree.
Its highest component has a primitive in the associated
graded cone.  Lift that primitive to the filtered cone and
subtract its differential.  The greatest degree decreases.
Induction terminates at degree zero and proves acyclicity.

Finally an acyclic cone makes $\chi$ a quasi-isomorphism.
Its two component equations give surjectivity on cohomology
when applied to a target cycle. Applying them to a source
cycle whose image is a boundary gives injectivity.
This proves the theorem without an unbounded convergence claim.
\end{proof}

\begin{example}[The odd self-bracket in the bar comparison]
Take $|x|=1$, $|t|=2$, $[x,x]=2t$, with $t$ central and
zero differential.  Then $sx$ is even and $st$ is odd.
The formulas give
\[
 \chi((sx)^2)=2[x|x],\qquad
 b_-\bigl(2[x|x]\bigr)=-2[t],\qquad
 D((sx)^2)=-2st.
\]
Thus the nonzero odd square appears in both chain complexes
with the same coefficient and sign.

For an abelian even generator $e$, the Lie chain coalgebra
is $\kk\oplus \kk(se)$, since $se$ is odd.
For an abelian odd generator $\theta$, it is instead
$\kk[s\theta]$, and
\[
 \chi((s\theta)^n)=n![\theta|\cdots|\theta].
\]
The factorial explains both the characteristic-zero
isomorphism and its behavior under weight completion.
\end{example}

The chain equivalence in Theorem~\ref{thm:olb-lie-bar} does
not make strict primitives invariant under a coalgebra
quasi-isomorphism.  For $U(\kk e)=\kk[e]$ with $e$ even, the strict
primitive complex of the bar contains $[e^m]$ for every
$m\geq1$. Its Lie chain source has primitive space $\kk(se)$.
The polynomial bar calculation above exhibits a class killed on passing
from strict primitive cohomology to full cohomology.

\subsection{Signed characters and primitive recognition}
\label{sec:olb-characters}

Give $e$ and $\theta$ from Section~\ref{sp:section} charge one,
independently of their cohomological degrees zero and one.
Charge adds under multiplication.
For a charge-graded super vector space with finite pieces,
its unsigned dimension series is
$\sum_n(\dim V_{n,\bar0}+\dim V_{n,\bar1})z^n$.
Its signed dimension series is
$\sum_n(\dim V_{n,\bar0}-\dim V_{n,\bar1})z^n$.
Both are formal power series.

The ordered basis of $\kk[e]\otimes\Lambda(\theta)$ gives
\[
 \sum_n\dim A_n\,z^n=\frac{1+z}{1-z},\qquad
 \sum_n(\dim A_{n,\bar0}-\dim A_{n,\bar1})z^n=1.
\]
For every $n>0$, the even basis vector is $e^n$ and the odd
basis vector is $e^{n-1}\theta$.
Give both generators primitive coproducts.
With zero differential, the primitive Lie algebra is
$\kk e\oplus \kk\theta$ by Theorem~\ref{sp:primitives}.
It is nonzero although the signed series equals that of $\kk$.

On the same algebra put $de=\theta$ and $d\theta=0$.
Then
\[
 d(e^n)=ne^{n-1}\theta,\qquad d(e^n\theta)=0.
\]
This preserves charge and the primitive coproduct.
Every positive-charge odd vector is a boundary, and no
positive-charge even vector is a cycle.
Thus its cohomology is $\kk$, and its primitive cohomology is
zero.  The underlying signed series has not changed.

The unsigned series likewise fails to determine the bracket, as the
Heisenberg and abelian algebras in Proposition~\ref{scalar-not-faithful}
show. Thus a character identity does not specify the bracket, the
differential, or a primitive recognition map.

\subsection{Finite charge bounds and completed bar maps}
\label{sec:olb-finite-completion}

Use the charge monoid $Q_+$ and the total-height tensor
$\bar\otimes_N$ of Section~\ref{sec:height-tensor}.
Proposition~\ref{monoidal-restriction} makes height restriction
symmetric monoidal. In the bar, every insertion and every cut preserves
total charge, so the equations of Proposition~\ref{prop:olb-bar}
restrict in these categories. All differentials must preserve charge.

Now let $\mathfrak g$ have positive $Q_+$-charges, a
charge-preserving differential, and charge-additive bracket.
Every homogeneous element of its augmentation ideal has positive height.
Consequently a bar word in charge $\beta$ has at most
$\htop(\beta)$ factors, and its total enveloping word filtration
is at most $\htop(\beta)$.
The same bound applies to $S(s\mathfrak g)_\beta$.
The maps $D$, $b_-$, $\chi$, and both coproducts
preserve charge.

For either chain coalgebra $C=\bigoplus_\beta C_\beta$,
define its completion by
\[
 (\widehat C)^n=\prod_\beta C_\beta^n,
 \qquad
 \widehat C=\bigoplus_n(\widehat C)^n,
\]
and define its tensor target by the same degreewise convention
on $\prod_{\beta,\gamma}(C_\beta\otimes C_\gamma)$.
Each internal tensor product still uses a direct sum of
cochain degrees.
All structure equations can be checked charge by charge,
with finitely many decompositions of each total charge.

\begin{corollary}\label{cor:olb-completed-lie-bar}
Under these positive-charge hypotheses, the Lie-chain map
extends to a coalgebra quasi-isomorphism
\[
 \widehat{(S(s\mathfrak g),D)}
 \xrightarrow{\ \widehat\chi\ }
 \widehat{B_-(U(\mathfrak g))}.
\]
It also restricts to every finite total-height tensor category.
\end{corollary}

\begin{proof}
All maps and filtrations in the proof of
Theorem~\ref{thm:olb-lie-bar} preserve charge.
Its quasi-isomorphism is therefore a quasi-isomorphism in
each charge, either by that proof componentwise or by taking
the summands of its cohomology isomorphism.

Proposition~\ref{cohomology-completion} shows that taking these
products preserves the quasi-isomorphism.
The coproduct identity holds in each pair of charges.
The proved monoidal restriction gives the finite-height
assertion with its specified tensor target.
\end{proof}

For a charge-one abelian odd generator $\theta$ of degree one,
$s\theta$ has degree zero.  The completed map sends the
allowed degree-zero series
\[
 \sum_{n\geq0}\frac{(s\theta)^n}{n!}
 \quad\longmapsto\quad
 \sum_{n\geq0}[\theta|\cdots|\theta].
\]
For an even generator $e$ of degree zero, the analogous bar
series has one term in every nonpositive cochain degree.
It is excluded by the direct sum over degrees.
Finite charge support, degreewise charge completion, and an
unrestricted product over word lengths are distinct choices.

In a relative bar construction with differential graded
coefficients, extension of coefficients is performed at each
specified finite weight before any completed product.
The ordinary tensor identity in
Section~\ref{sec:olb-relative-signs} does not assert that a
derived extension commutes with an unrestricted product.
A comparison over different coefficient algebras must therefore specify
the augmentations, base change and tensor targets.
An abstract enveloping realization supplies Lie and primitive
operations.  A geometric Hall realization must additionally
construct the geometric carrier and show that its product,
coproduct, differential and charge maps agree with them.
